\subsection{Certifying seat retention in Chile}\label{sec:retain-evidence}
For cohort $y$ and draw $s$, denote by $M^s$ the blind-menu allocation and by $N_D^s$ the allocation obtained by clearing a single policy $D$ shared across draws. Write $H(D)$ for its number of historical April hits. With $A$ denoting the blind policy, our finite-panel target for each cohort is $\max_D\{H(D)-H(A): X_s(D)=0\ \forall s\}$, where
\[
 X_s(D)=\sum_i \mathbf 1\{M^s(i)\ne\bot,\ N_D^s(i)=\bot\}.
\]
Seat retention requires $X_s(D)=0$ on every specified draw. This counts pupils across all grades, including applicants without an enrollment link. Net placement also counts entrants and can conceal an exit. A school switch requires the same pupil to receive seats under both menus and have different school identifiers. Multiple courses may belong to one school. Both prospective rules take the ministry's perspective after full reports, capacities and school-specific special-admission rankings have arrived. The ministry collects the special rankings separately. Neither rule consults the realized ordinary lottery order, initial offers or later cohort enrollment to choose a menu \citep{saeprioritytiming}. Our objective counts a linked pupil once if their historical April school is displayed. An applicant's historical enrollment is fixed for scoring.

\paragraph{An order-independent guarantee}
Call a course \emph{ample} if its total number of physical seats is at least the number of distinct pupils who could list it on either menu. All pupils may try all four types of physical seats in the reconstructed mechanism, with priority depending on the type. Even if a pupil is displaced from one type, they cannot exhaust an ample course. Otherwise every physical seat there would ultimately be occupied by a distinct other proposer who outranked them when they tried that category. A common continuity course also guarantees a seat. A third sufficient test certifies a specific common course for a pupil if, in any one category, fewer competitors can rank ahead of them than there are seats. For this count compare the pupil at their least favorable sibling priority with each possible rival at the rival's most favorable priority. Treat every unresolved tie as going to the rival.

Where a pupil has a guaranteed common option, courses listed after the first such option can be discarded for computing the outcome, except for any courses that could be promoted by an older relative. Also discard noncontinuity courses without physical seats. Treat an ample course as a terminal offer when computing course destinations. It can be treated in this way even if different pupils occupy different seat categories after the edit. To see this, hold fixed the set of pupils who ever reached that course in a DA run. A rejection at the last category would require at least as many \emph{other} distinct pupils at the course as there are seats. Since such a rejection is impossible, any further proposals following a displacement there stay within that course. Replacing these internal steps by a terminal offer changes no proposal at another course.

Consider all pupils lacking a guaranteed common option. Starting from them, include every potential active competitor at each positive-capacity nonample course on their effective lists. Include every older relative who could affect these pupils' preferences or priority. Repeat to closure. Freeze membership of every active course on the rows in this closure. Rows outside it may choose freely within the common quota and candidate pool. Every pupil outside has a guaranteed option. On the inside, all potential interactions at nonample courses and all relevant older relatives are included. Induction over grades and DA proposals after the terminal-offer replacement gives the same course destinations there on the old and new lists. Hence every baseline recipient stays placed, for every strict draw of the lottery. This proof permits an arbitrary change of seat type at an ample course and does not require any independence assumption on lottery orders. The prospective version fills available slots with the earliest submitted candidates. The hindsight version keeps every old hit and adds the April school wherever its row permits. Neither uses a realized lottery; only the hindsight version sees the April school.

Over the three cohorts this rule edits 6,186 pupils' menus. Ranking the permitted choices by submitted order gains 28 April hits in total ($-2,+5,+25$). Its hindsight version adds 118 ($29,32,57$) and loses no old hit. At least 118 gains are thus achievable even if the district requires seat retention under every lottery order admitted by the mechanism. In 96 independent held-out cohort--draws neither policy produced an incumbent exit. The rank-based rule generated 2,398 school switches across these simulations (the same pupil may switch on several draws). This universal certificate is conservative: only 99 changed rows in the rank-based menus contained edits on courses we could not exclude from a possible proposal path.

\paragraph{Constructing an upper and a lower bound on the same panel}
For the lower endpoint, start with pupils whose blind menu missed their April school. Add one course at that school and remove one old-only course, so every retained swap contributes one new hit and destroys none. Try the edits together, re-clear the entire national market, and withdraw edits associated with an exit on a design draw. We ran two searches. The first favors swaps compatible with the earlier exact-contract certificate. The second uses 64 design draws to avoid removing a course that often supplied the pupil's last attainable seat. In both searches the proposal histories determine which edits to reconsider. To certify a result, fix the final menus and re-clear them against the blind allocation on all 1,024 original draws, checking every pupil. The first search gained 625 hits across the cohorts. The second improved on it in all three, adding 170, 199 and 273 respectively. Rescoring the six final menu profiles directly against the registers confirmed every gain and zero prior hits lost. Across 3,072 complete national replays, the selected profiles had zero exits, 938 entrant events and 324,642 moves between courses among pupils who remained placed. These last two counts sum pupil--draws, including repeat observations of the same person. The searches may miss profitable combinations, including swaps that become safe only when taken together.

The starting upper bound is 4,823 pupils. Each one missed their April school under the blind menu and had a course at that school among the proposed additions. This is an upper bound on new hits even with unrestricted changes to all rows in the allowed family. One way to remove impossible gains is to find courses at which a pupil cannot secure a seat. For each draw, run DA on the longest common initial list of each pupil whose order and priority are unaffected by family assignments. Give the other pupils empty lists. Conditional on the eventual family assignments of any proposed menu, the retained reports are prefixes of the corresponding lists and their priorities agree. One can resume DA from the end of the prefix run, adding back omitted applicants and options below the prefixes. Physical seat pools can only become more selective. Thus a pupil who cannot attain any type at a course even under the prefix cutoffs and their most favorable possible sibling priority cannot end there. Retain the protected personal fallback. Enumerate the legal rows that show the pupil's April school. Eliminate a row if, on any draw where the blind menu places that pupil, no displayed course could seat them. A feasible policy's row must survive all such tests. Two further tests strengthen this bound. A pupil cannot end at a course if capacity-many other pupils must propose there first, lack a personal fallback there, and outrank them in every physical category. On each of the 1,024 draws per cohort we also ask whether all incumbents can fit simultaneously in the courses that survived the cutoff test. Start with the baseline assignment. Sending an incumbent to another course can preserve everyone only if a chain from that course reaches a vacancy or returns to their original course. This follows by decomposing any feasible change in assignments into paths and cycles. If a course must keep a particular set of incumbents to fill its physical seats, an outsider who outranks one of them in every category cannot propose there without displacing someone. For pupils with a fixed order, eliminate menus that must reach such a proposal. We intersect the surviving row configurations across these tests before counting possible gains. The cutoff test alone leaves 4,641 possible gains. Accounting for joint capacity and compulsory occupants rules out another 21, leaving 4,620. The first-proposal test rules out no further pupils after these exclusions. Surviving pupils may still compete for the same seat.

\paragraph{Learning from an earlier cohort}
Train the probability that a displayed candidate school is the historically attended school using the linked 2021 switchers. The ministry measured these labels in April 2022. Applications for the first evaluation cohort opened the following August. We assume a ministry-side implementation with access to the earlier administrative register and the school-specific priority inputs before assignment \citep{registerdata,saecalendar2022}. Group candidate schools by best submitted rank and by the maximum value of each of the two frozen offer scores for courses at that school. Shrink each of the 175 estimated cell probabilities toward its rank-group rate using 40 pseudo-observations, and shrink the rank group toward the overall rate in the same way. There are 82,467 candidate pupil--school pairs and 9,996 positive pairs in the calibration sample. The resulting probabilities are fixed in all three evaluation cohorts. They are predictions of a historical school label, not admission probabilities.

For each cohort first simulate the blind policy on the 1,024 design draws, then optimize the sum of predicted hits for distinct schools subject to the sufficient conditions below. Common courses and quotas stay fixed. Marginal credit for a second course at a school already shown is zero. The algorithm takes the highest-scoring distinct eligible schools in the remaining slots and fills any redundant slots, favoring the existing list at ties. This is rowwise optimal within the resulting restricted class. An oracle with the April register could gain exactly 163 hits within this class. The learned rule gains 50 ($9,22,19$), with 134 new hits and 84 previously covered pupils now missed. It changes 7,721 pupils' menus through 8,828 course substitutions. We froze the protocol, calibration coefficients and menu fingerprints before opening the test lotteries or scoring the menus on later registers. As an input check we altered all 1,580,726 recorded initial outcomes and reversed the ordinary tie priorities in the three archived markets, leaving school-supplied special rankings fixed. All six prospective menus retained their original hashes. Recomputing the learned rule's design restrictions on 12 sampled draws per year likewise returned identical restrictions.

\paragraph{A rule that lets incumbents move}
For each training lottery compute the blind allocation. In grades above pre-kindergarten, keep every old proposal up through the contract the pupil holds (all proposals for an unplaced pupil), and forbid adding a course that would be proposed to before that point. Intersect these tests across draws. The grade-by-grade proposal sequence, including family promotions and current-school fallbacks, is then identical in every higher grade.

For the last grade adopt a weaker test. Require the original course of every pupil placed on any training draw. A newly added course must either (i) fail to block the baseline in every training draw, or (ii) have at least as many ordinary and reserved seats in total as there are potential applicants in the union of the two policies. Hold common courses and individual quotas fixed. The added freedom in (ii) can change the allocated school.

\emph{Claim.} Every baseline recipient still has a seat on every training draw.
Fix one draw. The upper grades finish exactly as before, so the final grade faces unchanged contract priorities, preference orders and family states. For an incumbent assigned to a universally slack course (ii), truncate their new list after \emph{all} types at that course. For every other incumbent truncate immediately after the particular old contract, including its type. Previously unassigned pupils keep their whole new menus. At a non-(ii) pool any applicant proposing on the truncated lists either held a seat in that pool at baseline or ranks the contract above their old one. In the latter case stability on an old course, or the nonblocking test on a newly displayed course, means the pool was full with better-priority holders. This includes applicants who prefer a different type at their original course. The extra seat types allowed at the slack courses introduce no extra proposals at these ordinary pools. An old incumbent who reaches their floor pool is therefore among its highest-priority applicants and secures a seat. The guaranteed current-school fallback cannot reject its holder either. A pupil who reaches a universally slack course cannot be rejected by every type, since that would require capacity-many distinct \emph{other} applicants holding its seats. Thus at least one type on their truncated list accepts them. Every incumbent is therefore assigned on these truncated lists. Restoring options below each truncation cannot change the proposal path: no incumbent exhausted their list of courses in this run. The resulting full lists preserve all incumbent placements. A final-grade incumbent may switch schools, but their new course ranks at least as high in the institution's family-adjusted ordering. A pupil who stays at a slack course may receive a different seat type, and an initially unplaced pupil may receive a seat.


Re-clearing all 3,072 design markets produced zero exits, 829 entrants and 27,407 changes of school for the learned rule. On the 384 matched-course test draws it produced one exit (2023, draw 42), 104 entrants and 3,445 switches between schools. The unrestricted informed menus caused 277,663 exits over exactly the same draws. These totals count simulated pupil--draws. They do not estimate the number of actual children harmed. Within each cohort we treat independent simulated national draws as the units for a one-sided exact 95\% binomial upper limit on the probability of at least one exit under this generator. The limits for 2022--24 are 2.31\%, 3.65\%, and 2.31\%, respectively. They are marginal intervals. A second, pre-specified generator used a shared family ticket within schools. Across 32 fresh draws per year we observed one exit in 2022 and none in the later cohorts, 28 entrants and 879 school switches. The corresponding upper limits are 13.98\%, 8.94\%, and 8.94\%. We also applied the all-order rule on precisely these two panels, checking the baseline count on each paired draw. It produced 9,643 and 2,417 school switches, respectively, with zero exits or entrants in either panel. The learned rule had fewer switches yet failed to retain every incumbent in both tests. The identity of the person left without a seat matters even when total placement rises.

In a retrospective check of distributional incidence, pupils with the ministry's disadvantaged status accounted for 13 of the learned rule's 50 net April gains and 12 of the all-order rule's 28. Both rules gained coverage in that group in 2023 and 2024 and lost coverage in 2022. Among pupils with no original offer, neither rule gained or lost a single historical April hit in the linked sample. This stratification uses historical offers only after the menus are frozen. It does not identify a causal effect of recommendations.
